% GENERATED FILE. Edit canonical lessons, then run npm run generate:books.
% canonical-source-hash: fnv1a64:0d0aef474ae5c878
% canonical-lessons: UR-C251-guriya, UR-C251-phuldan, UR-C251-jinz, UR-C251-nekar, UR-C251-barsati

\chapter{Doll, Vase, Jeans, Shorts, Raincoat}
\label{ch:ur-doll-vase-jeans-shorts-raincoat}
\hlchaptermodality{251}

\begin{chapteropening}
\textbf{By the end of this chapter:} \emph{I can say a doll, a vase, jeans, shorts and a raincoat.}

Complete the last lesson of chapter 251: I can say a doll, a vase, jeans, shorts and a raincoat.
\end{chapteropening}

\section[\texorpdfstring{\ur{گڑیا} (guṛiyā)}{گڑیا (guṛiyā)}]{\ur{گڑیا} (guṛiyā) — a doll}

\label{lesson:UR-C251-guriya}



\begin{quote}
Before the new one: say the Urdu for a clay oven, then the Urdu for a kettle.
\end{quote}

\subsection*{You'll want to know: \ur{گڑیا}}
\textbf{\ur{گڑیا}} — \emph{guṛiyā} — \textquotedblleft{}a doll\textquotedblright{}.

\begin{cousinweb}
\textbf{\ur{گڑیا}} (\emph{guṛiyā}) is a tadbhava, a word inherited and reshaped by sound change. It comes from \emph{*guḍḍa}, \textquotedblleft{}doll, effigy\textquotedblright{}, a form CDIAL reconstructs from its descendants (CDIAL 4189).
\end{cousinweb}

\begin{grammarlens}[title={What you've built}]
One more word for this chapter.
\end{grammarlens}

\subsection*{Your turn}
\begin{itemize}
\raggedright
  \item \emph{Say it:} \emph{guṛiyā}
  \item \emph{Say it:} \emph{guṛiyā}, once more
  \item \emph{Say it:} say \emph{guṛiyā}
  \item \emph{Recall:} say the Urdu for a clay oven, then the Urdu for a kettle, then say \emph{guṛiyā} again
\end{itemize}

\begin{tcolorbox}[breakable,colback=teal!4,colframe=teal!35!black,title={Before you move on}]
What does \ur{گڑیا} mean? (\textquotedblleft{}A doll\textquotedblright{}.) Say it once more.
\end{tcolorbox}
\section[\texorpdfstring{\ur{پھولدان} (phūldān)}{پھولدان (phūldān)}]{\ur{پھولدان} (phūldān) — a vase}

\label{lesson:UR-C251-phuldan}



\begin{quote}
Before the new one: say the Urdu for a kettle, then the Urdu for a doll.
\end{quote}

\subsection*{You'll want to know: \ur{پھولدان}}
\textbf{\ur{پھولدان}} — \emph{phūldān} — \textquotedblleft{}a vase\textquotedblright{}.

\begin{cousinweb}
\textbf{\ur{پھولدان}} (\emph{phūldān}) is \textquotedblleft{}flower-holder\textquotedblright{}. \emph{phūl}, \textquotedblleft{}flower\textquotedblright{}, comes from Sanskrit \emph{phulla}, \textquotedblleft{}blown, opened out\textquotedblright{} (CDIAL 9092), and \emph{-dān}, \textquotedblleft{}holder, vessel\textquotedblright{}, is Persian.
\end{cousinweb}

\begin{grammarlens}[title={What you've built}]
One more word for this chapter.
\end{grammarlens}

\subsection*{Your turn}
\begin{itemize}
\raggedright
  \item \emph{Say it:} \emph{phūldān}
  \item \emph{Say it:} \emph{phūldān}, once more
  \item \emph{Say it:} say \emph{phūldān}
  \item \emph{Recall:} say the Urdu for a kettle, then the Urdu for a doll, then say \emph{phūldān} again
\end{itemize}

\begin{tcolorbox}[breakable,colback=teal!4,colframe=teal!35!black,title={Before you move on}]
What does \ur{پھولدان} mean? (\textquotedblleft{}A vase\textquotedblright{}.) Say it once more.
\end{tcolorbox}
\section[\texorpdfstring{\ur{جینز} (jīnz)}{جینز (jīnz)}]{\ur{جینز} (jīnz) — jeans}

\label{lesson:UR-C251-jinz}



\begin{quote}
Before the new one: say the Urdu for a doll, then the Urdu for a vase.
\end{quote}

\subsection*{You'll want to know: \ur{جینز}}
\textbf{\ur{جینز}} — \emph{jīnz} — \textquotedblleft{}jeans\textquotedblright{}.

\begin{cousinweb}
\textbf{\ur{جینز}} (\emph{jīnz}) is a loan from English \emph{jeans}.
\end{cousinweb}

\begin{grammarlens}[title={What you've built}]
One more word for this chapter.
\end{grammarlens}

\subsection*{Your turn}
\begin{itemize}
\raggedright
  \item \emph{Say it:} \emph{jīnz}
  \item \emph{Say it:} \emph{jīnz}, once more
  \item \emph{Say it:} say \emph{jīnz}
  \item \emph{Recall:} say the Urdu for a doll, then the Urdu for a vase, then say \emph{jīnz} again
\end{itemize}

\begin{tcolorbox}[breakable,colback=teal!4,colframe=teal!35!black,title={Before you move on}]
What does \ur{جینز} mean? (\textquotedblleft{}Jeans\textquotedblright{}.) Say it once more.
\end{tcolorbox}
\section[\texorpdfstring{\ur{نیکر} (nekar)}{نیکر (nekar)}]{\ur{نیکر} (nekar) — shorts}

\label{lesson:UR-C251-nekar}



\begin{quote}
Before the new one: say the Urdu for a vase, then the Urdu for jeans.
\end{quote}

\subsection*{You'll want to know: \ur{نیکر}}
\textbf{\ur{نیکر}} — \emph{nekar} — \textquotedblleft{}shorts\textquotedblright{}.

\begin{cousinweb}
\textbf{\ur{نیکر}} (\emph{nekar}) is a loan from English \emph{knickers}.
\end{cousinweb}

\begin{grammarlens}[title={What you've built}]
One more word for this chapter.
\end{grammarlens}

\subsection*{Your turn}
\begin{itemize}
\raggedright
  \item \emph{Say it:} \emph{nekar}
  \item \emph{Say it:} \emph{nekar}, once more
  \item \emph{Say it:} say \emph{nekar}
  \item \emph{Recall:} say the Urdu for a vase, then the Urdu for jeans, then say \emph{nekar} again
\end{itemize}

\begin{tcolorbox}[breakable,colback=teal!4,colframe=teal!35!black,title={Before you move on}]
What does \ur{نیکر} mean? (\textquotedblleft{}Shorts\textquotedblright{}.) Say it once more.
\end{tcolorbox}
\section[\texorpdfstring{\ur{برساتی} (barsātī)}{برساتی (barsātī)}]{\ur{برساتی} (barsātī) — a raincoat}

\label{lesson:UR-C251-barsati}



\begin{quote}
Before the new one: say the Urdu for jeans, then the Urdu for shorts.
\end{quote}

\subsection*{You'll want to know: \ur{برساتی}}
\textbf{\ur{برساتی}} — \emph{barsātī} — \textquotedblleft{}a raincoat\textquotedblright{}.

\begin{cousinweb}
\textbf{\ur{برساتی}} (\emph{barsātī}) is \textquotedblleft{}of the rains\textquotedblright{}: \emph{barsāt}, \textquotedblleft{}the rainy season\textquotedblright{}, from Sanskrit \emph{varṣārātri} (CDIAL 11398), with the ending \emph{-ī}.
\end{cousinweb}

\begin{grammarlens}[title={What you've built}]
One more word for this chapter.
\end{grammarlens}

\subsection*{Your turn}
\begin{itemize}
\raggedright
  \item \emph{Say it:} \emph{barsātī}
  \item \emph{Say it:} \emph{barsātī}, once more
  \item \emph{Say it:} say \emph{barsātī}
  \item \emph{Recall:} say the Urdu for jeans, then the Urdu for shorts, then say \emph{barsātī} again
\end{itemize}

\begin{tcolorbox}[breakable,colback=teal!4,colframe=teal!35!black,title={Before you move on}]
What does \ur{برساتی} mean? (\textquotedblleft{}A raincoat\textquotedblright{}.) Say it once more.
\end{tcolorbox}
